\section{Limitations}
\label{sec:limitations}

\paragraph{Theoretical scope at small $T$.}
Theorem~\ref{thm:forward-forgetting} is numerically vacuous for small $T$ at the
chosen hyperparameters ($\eps_1(T=20) \gg 1$). The bound becomes meaningful only
after a cooling-off period $T \gtrsim T^\star \approx 50$. For $T < T^\star$, the
feature-block contribution is loose, and the end-to-end guarantee is supplied by the
summary-subtraction cost $\mu_\summary$ via Lemma~\ref{lem:subtraction} and the
adaptive Gaussian-DP composition of Theorem~\ref{thm:composition}.

\paragraph{Non-convex formal scope.}
The rate $\mu_1(T)$ in Theorem~\ref{thm:forward-forgetting} is derived from the
linear structure of the aggregation state $\tilde\Delta$.
Appendix~\ref{app:nonconvex} supplies a trajectory-Lipschitz extension to
non-convex losses following \cite{gupta2021adaptive}, with the per-round Lipschitz
constant $L'$ measured empirically at ResNet-18 scale: $32$ (client, direction)
pairs give p95 $87$ and max $88$, which fall below the contraction threshold
$L' < 100$ imposed by the paper's hyperparameters. The trajectory-adjusted
sensitivity constant is therefore bounded by $\approx 9\,C$ at the operating
point, dominated by the $\alpha^T$ decay for any $T \gtrsim 20$. The random-direction
probe is a worst-case estimator, so this is an upper bound on the effective
non-convex correction factor at this operating point. A larger-scale probe
(ImageNet-class models, cross-task heterogeneity) and a formally sharp
non-convex bound are left to future work.

\paragraph{Formal-vs-empirical gap.}
Section~\ref{sec:results-tightness} applies a Gaussian-DP tightening of the
summary channel's MIA bound, which reduces the worst-case adversary advantage from
$0.76$ (standard $(\varepsilon, \delta)$-DP conversion) to $0.164$ at the operating
point $(\sigma_\summary \approx 2.42, \Delta = 1)$, i.e., a $4.6\times$ tightening.
The residual gap to the empirical LiRA of $0.046$ at the Carlini et al.\ reference
budget $n_\mathrm{shadow}{=}128$ (seed~0) is $0.12$. The main remaining source of
looseness is the head-only release being smaller than a worst-case Gaussian release
of the full summary. Appendix~\ref{app:rdp} supplies an additional R\'enyi-DP
tightening for the full-trajectory release model.

\paragraph{FedRecovery scoping.}
The destroy-at-every-finite-$\varepsilon$ finding in Section~\ref{sec:results} holds
for clip norms $C$ consistent with trainable FCL on ResNet-18. The reference
FedRecovery setting ($C \approx 1$) is incompatible with our training protocol:
measured local-update L2 norms are observed in $[6, 60]$, so $C = 1$ clips away the
majority of the gradient signal and training does not converge.

\paragraph{Training-substrate ambiguity.}
The retrain-from-scratch gold, the no-unlearning control, and the five baseline
methods are evaluated on a vanilla-FedAvg shared base; FORGE is evaluated on its own
decayed-aggregation substrate. A within-aggregator diagnostic isolates the two
contributions: FORGE trained with $\alpha{=}0$ and $\sigma{=}0$ (reducing the
decayed aggregator to vanilla FedAvg) reaches $0.1795$ accuracy at seed~$0$ with
retrain-KL $0.040$, versus $0.1763$ accuracy and retrain-KL $0.070$ for the main-text
decayed aggregator ($\alpha{=}0.5, \sigma{=}10^{-3}$) at the same seed. The
decayed-aggregator substrate therefore costs approximately $0.3$ pp of utility
relative to vanilla FedAvg, which accounts for roughly half of the $0.6$ pp gap to
SIFU in Table~\ref{tab:cifar10-ci-baselines}. The remaining $0.3$ pp is attributable to
the erasure step and training-run noise.

\paragraph{Scope of the class-incremental stress test.}
On the primary non-IID stationary benchmark (Table~\ref{tab:cifar10-benchmark})
retrain reaches $0.631 \pm 0.008$ accuracy, which is the standard FedAvg
operating point. The class-incremental stress test
(Table~\ref{tab:cifar10-ci-baselines}) runs the same protocol under 5 sequential
2-class tasks with no replay, where catastrophic forgetting pins preserved-utility
methods near $0.17$. The stress test is useful for a 12-method baseline
comparison (unlearning is hardest under catastrophic forgetting) but is not a
realistic deployment regime; we report both regimes to make the preserved-utility
ordering reproducible across protocols.

\paragraph{Adversarial-departure demonstration.}
The threat model (Section~\ref{sec:prelim}) admits a departing client that
contributed adversarially (e.g., a trigger-patch backdoor) before requesting
erasure. Lemma~\ref{lem:subtraction} removes the departing client's
head-space contribution independent of its input distribution, so the formal
guarantee carries through unchanged. An empirical attack-success-rate measurement
is not reported: at the paper's class-incremental hyperparameters (ResNet-18, no
replay, five sequential two-class tasks), the FCL protocol's catastrophic
forgetting already drives the baseline attack success rate to zero against any
single-client late-stage trigger before erasure, so an erasure-removal signal is
not observable at this operating point. A demonstration under a replay-enabled or IID-task protocol is
left to future work.

\paragraph{Summary-type ablation.}
The ablation across the two implemented summary types (moments and logits) is
reported in Section~\ref{sec:results-ablation} and Table~\ref{tab:ablation}. Two
alternative types in the released code (DP synthetic microbatches
\cite{harder2021dpmerf} and random-projection sketched gradients) are not evaluated
in this paper.

\paragraph{Summary-submission schedule at cross-device scale.}
The current implementation ties summary submission to the FL sampling schedule: a
client submits a summary only on rounds that coincide with its being sampled. At
the main-text operating point ($100$ clients, $20\%$ sample fraction, $F{=}20$),
each client submits approximately $2$ summaries per training run. At true
cross-device scale ($1000$ clients, $2\%$ sample fraction), the expected number of
submissions per client falls to approximately $0.2$, so most clients submit no
summary and cannot be served by the erasure procedure. The protocol-level
remediation is to decouple summary submission from the FL sampling schedule, i.e.,
require every client to submit at least one summary during a dedicated commitment
window at the start of training. This is a scheduling change rather than a change
to Algorithms~\ref{alg:decayed-aggregation}--\ref{alg:erasure}; the formal
guarantees carry through unchanged.
